\section*{\S\,3. Принцип двойственности}
\addcontentsline{toc}{section}{\S\,3. Принцип двойственности}
\markright{\S\,3. Принцип двойственности}
\subsection*{Условия}
1. Существует ли самодвойственная двухместная операция $F(A,B)$, не равная
тождественно $A$ или $B$?

2. Описать множество всех самодвойственных формул вида $F(A_1,\dots,A_n)$
при любом $n\ge2$.

3. Описать множество всех логических формул, для которых
$F\cap F^*=\emptyset$. Будет ли при этом $F\cup F^*=I$?

4.\fixed{в оригинале номер упражнения <<4.>> пропущен.} Будем
предполагать, что $F(A_1,\dots,A_n)$ "--- \emph{строго $n$-местная}
логическая операция, если не существует таких целых чисел
$k_1,\dots,k_s\in[1,n]$ ($s<n$), что
$F(A_1,\dots,A_n)=\Phi(A_{k_1},\dots,A_{k_s})$, где $\Phi$ "--- $s$-местная
логическая формула. Для каких $n$ существуют строго $n$-местные
самодвойственные логические формулы? Простой перебор показывает, что для
$n=2$ таких формул нет. Для $n=3$ примером самодвойственной строго
трехместной формулы является (1.6).

5. Пусть $\mu(n)$ "--- количество самодвойственных $n$-местных логических
формул. Показать, что $\mu(n)$ "--- неубывающая функция.

6. Пусть $\tau(n)$ "--- количество самодвойственных строго $n$-местных
логических формул. Показать, что $\tau(2)=0$. Для каких еще $n$
$\tau(n)=0$?

7. Пусть $\lambda\{F(A_1,\dots,A_n)\}=\lambda(F)$ "--- количество
конституентных областей в диаграмме Эйлера для логической формулы
$F(A_1,\dots,A_n)$, а $\lambda(F^*)$ "--- количество таких же областей для
двойственной формулы. Описать класс функций, для которых
$\lambda(F)=\lambda(F^*)$.

\subsection*{Решения}
% Решения упражнений § 3 (с. 21--22).
% Основа: заметки claude/p21-*.md из проекта (см. каталог notes/).

\begin{remarkbox}{Общие инструменты для упражнений 1--7}
Во всех упражнениях формулы рассматриваются \emph{с точностью до
тождественного равенства}: две формулы, тождественно равные друг другу,
считаются одной и той же (они задают одну и ту же <<операцию над
множествами>>).

\textbf{Лемма (L) о двойственности.} Для любой формулы $F$
\begin{equation*}\tag{L}
F^*(A_1,\dots,A_n)\equiv\ov{F(\ov{A_1},\dots,\ov{A_n})}.
\end{equation*}

\emph{Доказательство} "--- полная индукция по числу $k(F)$ символов
операций $\cap$, $\cup$, ${}^{\overline{\phantom a}}$ в записи $F$.
Обозначим $\widehat F=\ov{F(\ov A)}$, где $\ov A=(\ov{A_1},\dots,\ov{A_n})$.

\emph{База} $k=0$: $F=A_i$ дает $F^*=A_i=\ov{\ov{A_i}}=\widehat F$;
$F=\emptyset$ дает $F^*=I=\ov\emptyset=\widehat F$; $F=I$ дает
$F^*=\emptyset=\ov I=\widehat F$.

\emph{Шаг.} Пусть (L) верно для всех формул с не более чем $k$ операциями
и $k(F)=k+1$. Внешняя операция $F$ "--- одна из трех.
\begin{itemize}[nosep]
\item $F=\ov G$, $k(G)=k$. Замена символов черту не трогает, поэтому
  $F^*=\ov{G^*}\equiv\ov{\ov{G(\ov A)}}=G(\ov A)$, и
  $\widehat F=\ov{\ov{G(\ov A)}}=G(\ov A)$.
\item $F=G\cap H$, $k(G)+k(H)=k$, так что $k(G),k(H)\le k$ (именно здесь
  нужна \emph{полная} индукция: подформулы могут быть намного короче).
  Тогда $F^*=G^*\cup H^*\equiv\ov{G(\ov A)}\cup\ov{H(\ov A)}
  \overset{10^\circ}{\equiv}\ov{G(\ov A)\cap H(\ov A)}=\widehat F$.
\item $F=G\cup H$ "--- так же, с правилом $11^\circ$. \qquad$\square$
\end{itemize}

\textbf{Коды точек и булева функция формулы.} Точке $x\in I$ сопоставим
\emph{код} $\alpha(x)=(\alpha_1,\dots,\alpha_n)\in\{0,1\}^n$, $\alpha_i=1$
при $x\in A_i$ и $\alpha_i=0$ иначе. Точки с одним кодом $\alpha$ образуют
\emph{конституенту}
\[
K_\alpha=A_1^{\alpha_1}\cap\dots\cap A_n^{\alpha_n},\qquad A^1=A,\ A^0=\ov A ;
\]
конституент $2^n$, они попарно не пересекаются и покрывают $I$.
Принадлежность точки множеству $F(A_1,\dots,A_n)$ зависит только от ее
кода (индукция по $k(F)$), поэтому формуле отвечает булева функция
$f\colon\{0,1\}^n\to\{0,1\}$ (таблица истинности; $\cap\to\wedge$,
$\cup\to\vee$, черта $\to$ отрицание, $\emptyset\to0$, $I\to1$) и
\begin{equation*}\tag{D}
x\in F\iff f(\alpha(x))=1,\qquad F=\bigcup_{\alpha:\,f(\alpha)=1}K_\alpha .
\end{equation*}

\emph{(i) Тождества = равенство функций.} $F\equiv G\iff f=g$. Если
$f(\alpha_0)\ne g(\alpha_0)$, возьмем $I=\{0,1\}^n$,
$A_i=\{\alpha:\alpha_i=1\}$; тогда $K_{\alpha_0}=\{\alpha_0\}$, и точка
$\alpha_0$ лежит ровно в одном из множеств $F$, $G$.

\emph{(ii) Функция двойственной формулы.} Замена $A_i\to\ov{A_i}$
переводит код каждой точки в $\bar\alpha=(1-\alpha_1,\dots,1-\alpha_n)$.
Поэтому по (L)
\[
f^*(\alpha)=\ov{f(\bar\alpha)} .
\]
Наборы $\{0,1\}^n$ разбиваются на $2^{n-1}$ пар противоположных
$\{\alpha,\bar\alpha\}$ (набор не совпадает со своим дополнением).
\end{remarkbox}

%----------------------------------------------------------------------
\begin{solution}[1 (\S\,3)]
\emph{Ответ: формально да "--- это $\ov A$ и $\ov B$, но обе они
по существу одноместны; самодвойственной операции, существенно
зависящей от обоих аргументов, нет.}

По (L) самодвойственность $F\equiv F^*$ равносильна
\begin{equation*}\tag{C}
F(\ov{A_1},\dots,\ov{A_n})\equiv\ov{F(A_1,\dots,A_n)},
\qquad\text{т.\,е.}\qquad f(\bar\alpha)=\ov{f(\alpha)}\ \ \forall\alpha .
\end{equation*}
При $n=2$ пары противоположных наборов: $\{00,11\}$ и $\{01,10\}$.
Свободно выбираются $f(00)$ и $f(01)$, остальное определено:
\begin{center}\small
\begin{tabular}{ccccl}
\toprule
$f(00)$ & $f(01)$ & $f(10)$ & $f(11)$ & формула\\ \midrule
0&0&1&1& $A$\\ 0&1&0&1& $B$\\ 1&0&1&0& $\ov B$\\ 1&1&0&0& $\ov A$\\ \bottomrule
\end{tabular}
\end{center}
Всего четыре самодвойственные двухместные формулы, и ни одна не зависит
от обоих аргументов. Это и есть <<простой перебор>>, упомянутый в упр.~4.
\end{solution}

%----------------------------------------------------------------------
\begin{solution}[2 (\S\,3)]
\textbf{Разминка.}
\begin{center}\small
\begin{tabular}{lll}
\toprule
$F$ & $F^*$ & самодвойственна?\\ \midrule
$A_k$, $\ov{A_k}$ & те же & да\\
$\emptyset$ & $I$ & нет\\
$A\cap B$ & $A\cup B$ & нет (при $A=\emptyset$, $B=I$)\\
$(A\cap B)\cup(\ov A\cap I)$ & $(A\cup B)\cap(\ov A\cup\emptyset)$ & нет: $\ov A\cup B$ и $\ov A\cap B$\\
$(A\cap B)\cup(A\cap C)\cup(B\cap C)$ & $(A\cup B)\cap(A\cup C)\cap(B\cup C)$ & да "--- (1.6)\\
\bottomrule
\end{tabular}
\end{center}
(Проверка (L) на четвертой строке: при подстановке $A\to\ov A$ черта над
$A$ в формуле дает $\ov{\ov A}=A$; дополняются \emph{буквы}, а не
литералы: $F(\ov A,\ov B)=(\ov A\cap\ov B)\cup A=A\cup\ov B$,
$\ov{F(\ov A,\ov B)}=\ov A\cap B=F^*$.)

\textbf{Критерий.} По (C) и (ii) $F$ самодвойственна $\iff$
$f(\bar\alpha)=\ov{f(\alpha)}$ при всех $\alpha$. Наглядно: если
одновременно поменять ролями <<внутри>> и <<снаружи>> у \emph{всех}
исходных множеств, то результат тоже меняет <<внутри>> и <<снаружи>>.
Например, для (1.6) <<точка входит хотя бы в два из трех множеств>> после
дополнения всех трех превращается в <<входит не более чем в одно>>, т.\,е.
в дополнение исходного множества "--- проверка в одну строку.

На языке конституент: дополнение всех аргументов переставляет
$K_\alpha\leftrightarrow K_{\bar\alpha}$, и самодвойственность означает,
что формула из каждой пары $\{K_\alpha,K_{\bar\alpha}\}$ берет \emph{ровно
одну}.

\begin{theorem*}[описание]
$F(A_1,\dots,A_n)$ самодвойственна тогда и только тогда, когда
\[
F\equiv\bigcup_{\alpha\in S}A_1^{\alpha_1}\cap\dots\cap A_n^{\alpha_n},
\]
где $S\subset\{0,1\}^n$ содержит ровно один элемент из каждой пары
$\{\alpha,\bar\alpha\}$. Разным $S$ отвечают разные формулы.
\end{theorem*}
\begin{proof}
<<Только тогда>>: $S=\{\alpha:f(\alpha)=1\}$ и (C). <<Тогда>>:
$f(\alpha)=1\iff\alpha\in S\iff\bar\alpha\notin S\iff f(\bar\alpha)=0$.
Различность "--- по (i).
\end{proof}
Следствия: самодвойственная формула занимает ровно половину диаграммы
Эйлера, $\lambda(F)=2^{n-1}$; их число $\mu(n)=2^{2^{n-1}}$.

\begin{theorem*}[порождение]
$F$ самодвойственна тогда и только тогда, когда
\[
\boxed{\;F\equiv\bigl(A_1\cap G(A_2,\dots,A_n)\bigr)\cup\bigl(\ov{A_1}\cap G^*(A_2,\dots,A_n)\bigr)\;}
\]
для некоторой \emph{произвольной} формулы $G$ от $n-1$ аргументов;
соответствие $G\mapsto F$ взаимно однозначно.
\end{theorem*}
\begin{proof}
Для любой формулы верно разложение Шеннона
$F\equiv(A_1\cap F_1)\cup(\ov{A_1}\cap F_0)$, $F_1=F(I,A_2,\dots)$,
$F_0=F(\emptyset,A_2,\dots)$. Подставим его в (C). Слева
$(\ov{A_1}\cap F_1(\ov{A'}))\cup(A_1\cap F_0(\ov{A'}))$, где
$A'=(A_2,\dots,A_n)$; справа
$(A_1\cap\ov{F_1})\cup(\ov{A_1}\cap\ov{F_0})$. Множество вида
$(A_1\cap X)\cup(\ov{A_1}\cap Y)$ однозначно определяет $X$ и $Y$ (при
$A_1=I$ и $A_1=\emptyset$), поэтому
$F_1(\ov{A'})=\ov{F_0(A')}$, т.\,е. $F_0=\ov{F_1(\ov{A'})}\overset{(L)}{=}F_1^*$.
Обратно, для $F$ указанного вида по (L) $G(\ov{A'})=\ov{G^*(A')}$,
$G^*(\ov{A'})=\ov{G(A')}$, и
$F(\ov A)=(\ov{A_1}\cap\ov{G^*})\cup(A_1\cap\ov G)=\ov{F(A)}$.
Наконец, $G=F(I,A_2,\dots,A_n)$ восстанавливается по $F$.
\end{proof}
Это второй вывод $\mu(n)=2^{2^{n-1}}$: самодвойственных формул от $n$
аргументов столько же, сколько вообще формул от $n-1$ аргумента.

\textbf{Проверка при $n=3$.} Перебирая 16 формул $G(B,C)$, получаем
$F=(A\cap G)\cup(\ov A\cap G^*)$:
\begin{center}\small
\begin{tabular}{lll}
\toprule
$G$ & $G^*$ & $F$\\ \midrule
$\emptyset$ & $I$ & $\ov A$\\
$I$ & $\emptyset$ & $A$\\
$B$, $\ov B$, $C$, $\ov C$ & те же & $B$, $\ov B$, $C$, $\ov C$\\
$B^b\cap C^c$ & $B^b\cup C^c$ & $\maj(\ov A,B^b,C^c)$\\
$B^b\cup C^c$ & $B^b\cap C^c$ & $\maj(A,B^b,C^c)$\\
$B\sd C$ & $\ov{B\sd C}$ & $\ov{A\sd B\sd C}$\\
$\ov{B\sd C}$ & $B\sd C$ & $A\sd B\sd C$\\ \bottomrule
\end{tabular}
\end{center}
Здесь $b,c\in\{0,1\}$, $\maj(X,Y,Z)=(X\cap Y)\cup(X\cap Z)\cup(Y\cap Z)$.
Итого 16 формул: 6 вырожденных ($A,\ov A,B,\ov B,C,\ov C$), 8
<<медиан>> $\maj(A^a,B^b,C^c)$ (среди них (1.6)) и 2 <<четностных>>
$A\sd B\sd C$, $\ov{A\sd B\sd C}$.

\textbf{Итог.} Самодвойственные формулы "--- в точности те, что
дополняют результат при дополнении всех аргументов:
$F(\ov{A_1},\dots,\ov{A_n})\equiv\ov{F(A_1,\dots,A_n)}$. Их можно выписать
(1) через конституенты, выбирая по одной из каждой пары
$\{K_\alpha,K_{\bar\alpha}\}$, или (2) как
$(A_1\cap G)\cup(\ov{A_1}\cap G^*)$ с произвольной $G$.

\textbf{Замечания.} (а) \emph{Замкнутость.} Если $F$ и $G_1,\dots,G_n$
самодвойственны, то и $F(G_1,\dots,G_n)$ самодвойственна; вместе
с проекциями $A_k$ это замкнутый класс (класс $S$ в теории Поста) "---
ср. \S\,4. (б) \emph{Связь с $R$-функциями.} При переходе к функциям
($A_i\leftrightarrow\{\omega_i\ge0\}$) дополнение множества "--- это смена
знака $\omega\mapsto-\omega$, и (C) превращается в
$f(-\omega_1,\dots,-\omega_n)=-f(\omega_1,\dots,\omega_n)$:
самодвойственным операциям отвечают \emph{нечетные} $R$-функции.
Например, $\tfrac12(x+y\mp|x-y|)$ ($\min$ и $\max$) двойственны, но не
самодвойственны, а медиана $\operatorname{med}(x,y,z)$ нечетна, как
и должно быть для (1.6).
\end{solution}

%----------------------------------------------------------------------
\begin{solution}[3 (\S\,3)]
\emph{Ответ: $F\cap F^*\equiv\emptyset$ $\iff$
$F(\ov{A_1},\dots,\ov{A_n})\equiv F(A_1,\dots,A_n)$ (<<четные>> формулы);
для них всегда $F^*\equiv\ov F$ и потому $F\cup F^*\equiv I$; более того,
условие $F\cup F^*\equiv I$ выделяет тот же класс.}

\emph{Перевод на язык функций.} По (i) $F\cap F^*\equiv\emptyset$
$\iff$ $f(\alpha)\wedge f^*(\alpha)=0$ при всех $\alpha$, а по (ii)
это $f(\alpha)\wedge\ov{f(\bar\alpha)}=0$:
\begin{equation*}\tag{1}
f(\alpha)=1\ \Rightarrow\ f(\bar\alpha)=1\qquad\forall\alpha .
\end{equation*}
Применяя (1) к $\bar\alpha$ вместо $\alpha$ (а $\bar{\bar\alpha}=\alpha$),
получаем обратную импликацию, т.\,е.
\begin{equation*}\tag{2}
f(\alpha)=f(\bar\alpha)\qquad\forall\alpha,
\qquad\text{или}\qquad F(\ov{A_1},\dots,\ov{A_n})\equiv F(A_1,\dots,A_n).
\end{equation*}
Обратно, (2) $\Rightarrow$ (1) очевидно. То же в множествах: $F\cap F^*=\emptyset$
значит $F^*\subseteq\ov F$, т.\,е. (по (L)) $\ov{F(\ov A)}\subseteq\ov{F(A)}$,
$F(A)\subseteq F(\ov A)$ при любых $A$; подставив $\ov A$ вместо $A$,
получаем обратное включение.

\begin{theorem*}
$F\cap F^*\equiv\emptyset$ тогда и только тогда, когда $F$ "---
объединение \emph{целых пар} конституент:
$F\equiv\bigcup_{\alpha\in T}(K_\alpha\cup K_{\bar\alpha})$, где $T$ "---
произвольное множество представителей пар. Таких формул $2^{2^{n-1}}$.
\end{theorem*}

\emph{Замкнутая форма.} При $n\ge2$ четные формулы "--- это в точности
\[
\boxed{\;F\equiv H\bigl(A_1\sd A_2,\ A_1\sd A_3,\ \dots,\ A_1\sd A_n\bigr)\;}
\]
с \emph{произвольной} формулой $H$ от $n-1$ аргумента. Действительно,
$\ov{A_1}\sd\ov{A_i}\equiv A_1\sd A_i$, так что такая $F$ четна. Обратно,
для четной $F$ положим $H(B_2,\dots,B_n)=F(\emptyset,B_2,\dots,B_n)$. Если
код точки $x$ есть $\alpha$ и $\alpha_1=0$, то $x\in A_1\sd A_i\iff x\in A_i$,
и правая часть содержит $x$ тогда же, когда $f(\alpha)=1$; если
$\alpha_1=1$, то $x\in A_1\sd A_i\iff x\notin A_i$, код равен
$(0,\bar\alpha_2,\dots,\bar\alpha_n)=\bar\alpha$, и по (2)
$f(\bar\alpha)=f(\alpha)$. Неформально: четная формула <<не видит>> самих
множеств "--- только то, \emph{совпадают или различаются} состояния точки
по отношению к разным множествам.

\emph{Второй вопрос.} Для четной $F$ по (ii) $f^*(\alpha)=\ov{f(\bar\alpha)}=\ov{f(\alpha)}$,
т.\,е. $F^*\equiv\ov F$; значит $F\cup F^*\equiv I$ и $F\cap F^*\equiv\emptyset$
одновременно. Обратно, $F\cup F^*\equiv I$ означает
$f(\alpha)\vee\ov{f(\bar\alpha)}=1$, т.\,е. $f(\bar\alpha)=1\Rightarrow f(\alpha)=1$,
и тем же приемом получаем (2). Итак, для любой $F$ равносильны:
\[
F\cap F^*\equiv\emptyset\iff F\cup F^*\equiv I\iff F^*\equiv\ov F\iff
F(\ov{A_1},\dots,\ov{A_n})\equiv F(A_1,\dots,A_n).
\]
Для произвольной формулы ни одно из первых двух условий не влечет
другое <<по общим соображениям>>: например, для $F=A\cap B$ имеем
$F\cap F^*=A\cap B\ne\emptyset$ и $F\cup F^*=A\cup B\ne I$.

\textbf{Примеры.} $n=1$: четные формулы $\emptyset$ и $I$ ($A$ и $\ov A$
не подходят: $A\cap A^*=A$). $n=2$: $\emptyset$, $A\sd B$,
$\ov{A\sd B}=(A\cap B)\cup(\ov A\cap\ov B)$, $I$ (т.\,е. $H=\emptyset$,
$B_2$, $\ov{B_2}$, $I$); например,
$(A\sd B)^*=(A\cup\ov B)\cap(\ov A\cup B)=\ov{A\sd B}$. $n=3$: 16 формул
$H(A\sd B,A\sd C)$; при $H=B_2\cap B_3$ получаем
$F=(A\sd B)\cap(A\sd C)=K_{100}\cup K_{011}$, и
$F^*=(A\cup\ov B\cup\ov C)\cap(\ov A\cup B\cup C)=\ov F$.

\textbf{Связь с упр.~2.} Самодвойственные формулы ($f(\bar\alpha)=\ov{f(\alpha)}$,
<<нечетные>>) и четные ($f(\bar\alpha)=f(\alpha)$) "--- два крайних
случая: у первых $F\cap F^*=F$, у вторых $F\cap F^*=\emptyset$. Классы
равномощны, и $F$ самодвойственна $\iff$ $F\sd A_1$ четна.
\end{solution}

%----------------------------------------------------------------------
\begin{solution}[4 (\S\,3)]
\emph{Ответ: строго $n$-местные самодвойственные формулы существуют при
$n=1$ и при всех $n\ge3$; единственное исключение "--- $n=2$.}

\emph{Переформулировка.} Будем говорить, что $F$ \emph{существенно
зависит} от $A_j$, если найдутся $A_1,\dots,A_n$ и $A_j'$ с
$F(\dots,A_j,\dots)\ne F(\dots,A_j',\dots)$. Строгая $n$-местность "---
это существенная зависимость от каждого аргумента: если $F$ не зависит
от $A_j$, то $F\equiv F(A_1,\dots,\emptyset,\dots,A_n)$ "--- формула от
$n-1$ аргумента; и наоборот.

\emph{$n=2$} "--- см. решение упр.~1: все четыре самодвойственные формулы
одноместны. (Самодвойственных формул от двух аргументов $2^{2^1}=4$, ровно
столько же, сколько одноместных самодвойственных $A,\ov A,B,\ov B$ "---
<<места не остается>>.)

\emph{$n=3$}: формула (1.6) $M(A,B,C)$ самодвойственна; при $B=I$,
$C=\emptyset$ она равна $A$, и по симметрии существенно зависит от
каждого аргумента.

\emph{Общий случай $n\ge3$.} По теореме о порождении (упр.~2) при
$G=A_2\cap\dots\cap A_n$, $G^*=A_2\cup\dots\cup A_n$ получаем
самодвойственную формулу
\begin{equation*}\tag{4}
F_n=(A_1\cap A_2\cap\dots\cap A_n)\cup\bigl(\ov{A_1}\cap(A_2\cup\dots\cup A_n)\bigr)
\end{equation*}
(<<точка лежит в $A_1$ и во всех остальных множествах, либо не лежит в
$A_1$, но лежит хотя бы в одном из остальных>>; при $n=3$ это
$M(\ov A,B,C)$). Прямая проверка: при $P=A_2\cap\dots\cap A_n\subseteq U=A_2\cup\dots\cup A_n$
\[
F_n^*=(A_1\cup U)\cap(\ov{A_1}\cup P)=(A_1\cap P)\cup(\ov{A_1}\cap U)\cup(U\cap P)
=(A_1\cap P)\cup(\ov{A_1}\cap U)=F_n,
\]
так как $U\cap P=P=(A_1\cap P)\cup(\ov{A_1}\cap P)$ поглощается.
Существенная зависимость: при $A_2=I$, $A_3=\dots=A_n=\emptyset$ (нужен
аргумент $A_3$, т.\,е. $n\ge3$) $F_n=\ov{A_1}$; при $A_1=I$ и всех
$A_i=I$, кроме $A_j$ ($j\ge2$), $F_n=A_j$. При $n=2$ та же формула
вырождается в $(A_1\cap A_2)\cup(\ov{A_1}\cap A_2)=A_2$.

\emph{$n=1$}: $A$ (и $\ov A$) самодвойственна и не равна тождественно ни
$\emptyset$, ни $I$.

\emph{Оценка подсчетом.} Нестрогая самодвойственная формула равна самодвойственной формуле от
$n-1$ аргумента, поэтому их не больше $n\,\mu(n-1)$, тогда как
$\mu(n)=\mu(n-1)^2$. При $n=3$: $16>3\cdot4=12$, при $n=4$: $256>64$;
при $n=2$ оценка дает $2\mu(1)=4=\mu(2)$ "--- запаса нет, и перебор
показывает, что строго двухместных действительно нет.
\end{solution}

%----------------------------------------------------------------------
\begin{solution}[5 (\S\,3)]
\emph{Что считается.} Различных \emph{записей} бесконечно много
($A_1$, $A_1\cap A_1$, \dots), но формул с точностью до тождества
$2^{2^n}$ "--- по числу таблиц $f$ (см. (i)). Самодвойственные среди них
и составляют $\mu(n)$. При этом <<$n$-местная>> формула не обязана
существенно зависеть от всех аргументов.

\emph{Монотонность без вычислений.} Всякая самодвойственная формула
$F(A_1,\dots,A_n)$ является самодвойственной и как формула от
$A_1,\dots,A_n,A_{n+1}$ (дополнение лишнего аргумента ни на что не
влияет), а различные формулы остаются различными. Значит,
$\mu(n)\le\mu(n+1)$.

\emph{Точное значение.} По критерию (C) таблица задается выбором значения
на одном наборе каждой из $2^{n-1}$ пар $\{\alpha,\bar\alpha\}$, выборы
независимы:
\[
\mu(n)=2^{2^{n-1}}:\qquad
\begin{array}{c|cccccc}
n&1&2&3&4&5&6\\\hline
\mu(n)&2&4&16&256&65\,536&4\,294\,967\,296
\end{array}
\]
Отсюда $\mu(n+1)=\mu(n)^2>\mu(n)$ "--- функция даже строго возрастает.
Для сравнения: всех $n$-местных формул $2^{2^n}=\mu(n)^2=\mu(n+1)$,
т.\,е. самодвойственных формул от $n$ аргументов ровно столько, сколько
вообще формул от $n-1$ аргумента.
\end{solution}

%----------------------------------------------------------------------
\begin{solution}[6 (\S\,3)]
\emph{Ответ: $\tau(2)=0$, и ни при каком другом $n\ge1$ $\tau(n)$ не
обращается в нуль: $\tau(1)=2$, $\tau(3)=10$, $\tau(4)=208$,
$\tau(5)=64\,386$, \dots}

$\tau(2)=0$ "--- решение упр.~1; $\tau(1)\ge1$ ($A$) и $\tau(n)\ge1$ при
$n\ge3$ (формула (4) из решения упр.~4).

\emph{Точная формула.} У каждой формулы есть однозначно определенное
множество $S\subseteq\{1,\dots,n\}$ существенных аргументов; подставив
$\emptyset$ вместо несущественных, получаем строго $|S|$-местную формулу
$\Phi$, которая самодвойственна вместе с $F$. Обратно, пара <<$S$ +
строго $|S|$-местная самодвойственная $\Phi$>> дает самодвойственную
$n$-местную формулу, и разные пары дают разные формулы. Поэтому
\[
\mu(n)=\sum_{k=0}^{n}\binom nk\tau(k),\qquad \tau(0)=0
\]
(формулы $\emptyset$, $I$ от нуля аргументов не самодвойственны). Отсюда
$\tau(n)=\mu(n)-\sum_{k<n}\binom nk\tau(k)$ (или, обращением,
$\tau(n)=\sum_k(-1)^{n-k}\binom nk\mu(k)$):
\[
\begin{aligned}
\tau(1)&=2,\qquad \tau(2)=4-2\cdot2=0,\qquad \tau(3)=16-3\cdot0-3\cdot2=10,\\
\tau(4)&=256-4\cdot10-6\cdot0-4\cdot2=208,\\
\tau(5)&=65\,536-5\cdot208-10\cdot10-10\cdot0-5\cdot2=64\,386 .
\end{aligned}
\]
(Для $n\le4$ значения проверены прямым перебором.) При $n=3$ десять
формул "--- восемь <<медиан>> $\maj(A^{\pm},B^{\pm},C^{\pm})$ и две
<<четности>>.

\emph{Почему именно $n=2$.} Вырожденных самодвойственных формул не
больше $n\,\mu(n-1)$, а $\mu(n)/(n\,\mu(n-1))=\mu(n-1)/n=2^{2^{n-2}}/n$.
При $n\ge3$ это больше единицы и быстро растет, а при $n=2$ равно
$\mu(1)/2=1$: двойка "--- единственное $n$, при котором возведение
в квадрат еще не обгоняет умножение на $n$ ($2^2=2\cdot2$). Доля строго
$n$-местных среди самодвойственных быстро стремится к единице:
$\tau(4)/\mu(4)=208/256$, $\tau(5)/\mu(5)\approx0{,}98$.
\end{solution}

%----------------------------------------------------------------------
\begin{solution}[7 (\S\,3)]
\emph{Ответ: $\lambda(F)=\lambda(F^*)$ $\iff$ $\lambda(F)=2^{n-1}$, т.\,е.
$F$ составлена ровно из половины конституентных областей (равносильно:
$\lambda(F)=\lambda(\ov F)$). Таких формул $\binom{2^n}{2^{n-1}}$.}

$\lambda(F)$ "--- число конституент, из которых состоит $F$ на полной
диаграмме Эйлера (все $2^n$ областей присутствуют), т.\,е. число единиц
в таблице $f$. Например, при $n=2$: $\lambda(A\cap B)=1$,
$\lambda(A\cup B)=3$, $\lambda(A)=2$, $\lambda(\emptyset)=0$,
$\lambda(I)=4$.

По (ii) $f^*(\alpha)=1\iff f(\bar\alpha)=0$: области $F^*$ "--- это
наборы, противоположные наборам, \emph{не} вошедшим в $F$. Отображение
$\alpha\mapsto\bar\alpha$ "--- перестановка, поэтому
\[
\boxed{\;\lambda(F)+\lambda(F^*)=2^n\;}
\]
(геометрически: $F^*$ "--- центральное отражение дополнения $\ov F$ на кубе
кодов). Значит, $\lambda(F)=\lambda(F^*)\iff\lambda(F)=2^{n-1}$; назовем
такие формулы \emph{уравновешенными}. Любое множество $S$ из $2^{n-1}$
наборов дает уравновешенную формулу, разные $S$ "--- разные формулы:
\[
\binom{2^n}{2^{n-1}}:\quad n=1:\ 2,\quad n=2:\ 6,\quad n=3:\ 70,\quad n=4:\ 12\,870,\ \dots
\]

\textbf{Примеры.} $n=2$: $A$, $\ov A$, $B$, $\ov B$, $A\sd B=K_{01}\cup K_{10}$,
$\ov{A\sd B}=K_{00}\cup K_{11}$. $n=3$: медиана (1.6) "--- 4 области из 8;
$A\sd B\sd C$ "--- 4 области; $A\cap(B\cup C)$ "--- 3 области "--- \emph{не}
уравновешена, ее двойственная $A\cup(B\cap C)$ занимает 5.

\textbf{Связь с упр.~2 и 3.} Самодвойственные формулы берут ровно одну
область из каждой пары $\{K_\alpha,K_{\bar\alpha}\}$, поэтому \emph{все}
уравновешены. Четные берут пары целиком, $m$ пар дают $\lambda=2m$
($0\le m\le2^{n-1}$; крайние случаи $\emptyset$ и $I$); уравновешены
лишь те, что составлены ровно из половины пар, $m=2^{n-2}$, их
$\binom{2^{n-1}}{2^{n-2}}$ (при $n=2$ "--- $A\sd B$ и $\ov{A\sd B}$,
но не $\emptyset$, $I$; при $n=3$ "--- 6 из 16). Уравновешенная формула не
обязана быть ни той, ни другой: при $n=3$, $S=\{000,111,001,010\}$,
\[
F=K_{000}\cup K_{111}\cup K_{001}\cup K_{010}=\ov{A\sd(B\cap C)}
\]
берет пару $\{000,111\}$ целиком (не самодвойственна), а пару
$\{001,110\}$ "--- наполовину (не четна), но $\lambda(F)=4=\lambda(F^*)$.
При $n=3$ из 70 уравновешенных формул 16 самодвойственных, 6 четных
и 48 <<прочих>>.
\end{solution}
